\documentclass[12pt]{ximera}
\title{Practice Exam 2}
\begin{document}
\begin{abstract}
Practice for Exam 2, covering Chapter 2: critical players and veto power, counting
coalitions and sequential coalitions, Banzhaf and Shapley-Shubik power, and a city
council analyzed with coalition types.
\end{abstract}
\maketitle

\section*{Practice Exam 2}

\begin{problem}
Assess whether each of the following claims is true or false.

\begin{prompt}
\textbf{(a)} Claim: in a weighted voting system, it is possible for every player in a
\emph{winning} coalition to be critical.

\begin{multipleChoice}
\choice[correct]{True}
\choice{False}
\end{multipleChoice}

\begin{feedback}
True. In a winning coalition that only just reaches the quota, losing any member drops
it below. For example, in $[12:8,6,4,2]$ the coalition $\set{B,C,D}$ has weight exactly
$12$, so all three members are critical.
\end{feedback}
\end{prompt}

\begin{prompt}
\textbf{(b)} Claim: in a weighted voting system, it is possible for a player with weight
less than half of the quota to have veto power.

\begin{multipleChoice}
\choice[correct]{True}
\choice{False}
\end{multipleChoice}

\begin{feedback}
True. Veto power only requires that the other players together fall short of the quota.
In $[10:4,3,2,1]$ the quota equals the total weight, so every player --- including the
one with weight $1$, far below half of $10$ --- can block any motion.
\end{feedback}
\end{prompt}

\begin{prompt}
\textbf{(c)} Claim: in a weighted voting system, if two players have the same number of
votes, then they have equal power in the system.

\begin{multipleChoice}
\choice[correct]{True}
\choice{False}
\end{multipleChoice}

\begin{feedback}
True. Two players with equal weight are interchangeable: swapping them turns every
winning coalition into another winning coalition, so each is critical (or pivotal)
exactly as often as the other.
\end{feedback}
\end{prompt}
\end{problem}

\begin{problem}
Use counting of subsets and sequences to answer the following.

\begin{prompt}
\textbf{(a)} In a weighted voting system with $10$ players, how many coalitions have two
or more players?

\begin{multipleChoice}
\choice{$2^{10}=1024$}
\choice{$2^{9}=512$}
\choice[correct]{$2^{10}-10-1=1013$}
\choice{$2^{10}-1=1023$}
\end{multipleChoice}

\begin{feedback}
Of the $2^{10}=1024$ subsets, $1$ is empty and $10$ contain a single player. Removing
those leaves $1024-10-1=1013$.
\end{feedback}
\end{prompt}

\begin{prompt}
\textbf{(b)} In a system with $15$ players, how many coalitions include $P_1$ and exclude
$P_{12}$, $P_{13}$, $P_{14}$, and $P_{15}$?

$\answer{1024}$

\begin{feedback}
Five players are settled, so $10$ remain free to be in or out: $2^{10}=1024$.
\end{feedback}
\end{prompt}

\begin{prompt}
\textbf{(c)} In the system $[100:60,51,45,37,26,18]$ with six players, how many
\emph{sequential} coalitions are there?

$\answer{720}$

\begin{feedback}
A sequential coalition orders all six players: $6!=720$. The weights do not matter for
this count.
\end{feedback}
\end{prompt}
\end{problem}

\begin{problem}
Consider the weighted voting system $[12:8,6,4,2]$ with players A, B, C, and D.

\begin{prompt}
\textbf{(a)} How many winning coalitions are there?

$\answer{7}$

\begin{feedback}
The coalitions with weight at least $12$ are
\[
\set{A,B}=14,\ \set{A,C}=12,\ \set{A,B,C}=18,\ \set{A,B,D}=16,\ \set{A,C,D}=14,\
\set{B,C,D}=12,\ \set{A,B,C,D}=20.
\]
Critical players: both in $\set{A,B}$ and $\set{A,C}$; only A in $\set{A,B,C}$; A and B
in $\set{A,B,D}$; A and C in $\set{A,C,D}$; all three in $\set{B,C,D}$; none in the
grand coalition.
\end{feedback}
\end{prompt}

\begin{prompt}
\textbf{(b)} Claim: player D is critical in the winning coalition $\set{B,C,D}$.

\begin{multipleChoice}
\choice[correct]{True}
\choice{False}
\end{multipleChoice}

\begin{feedback}
True. $\set{B,C,D}$ has weight $6+4+2=12$; without D it drops to $10<12$.
\end{feedback}
\end{prompt}

\begin{prompt}
\textbf{(c)} Claim: this system is \emph{equivalent} to $[4:3,2,1,0]$.

\begin{multipleChoice}
\choice{True}
\choice[correct]{False}
\end{multipleChoice}

\begin{feedback}
False. In $[4:3,2,1,0]$ the coalition $\set{B,C,D}$ has weight $2+1+0=3<4$ and loses,
but in $[12:8,6,4,2]$ it wins. Different winning coalitions means the systems are not
equivalent.
\end{feedback}
\end{prompt}

\begin{prompt}
\textbf{(d)} Determine the Banzhaf power distribution, as percentages rounded to two
decimal places.

A: $\answer[tolerance=0.005]{41.67}\%$ \quad B: $\answer[tolerance=0.005]{25}\%$ \quad
C: $\answer[tolerance=0.005]{25}\%$ \quad D: $\answer[tolerance=0.005]{8.33}\%$

\begin{feedback}
Critical counts from part (a): A $5$, B $3$, C $3$, D $1$, total $12$. So A
$\frac{5}{12}\approx 41.67\%$, B and C $\frac{3}{12}=25\%$ each, D $\frac1{12}\approx 8.33\%$.
\end{feedback}
\end{prompt}
\end{problem}

\begin{problem}
The following are all $6$ sequential coalitions for a weighted voting system with three
players A, B, and C, with the pivotal player underlined.
\[
\begin{array}{ccc}
\seq{A,B,\underline{C}} & \seq{B,A,\underline{C}} & \seq{C,\underline{A},B}\\
\seq{A,\underline{C},B} & \seq{B,\underline{C},A} & \seq{C,\underline{B},A}
\end{array}
\]

\begin{prompt}
\textbf{(a)} Determine the Shapley-Shubik power distribution, as fractions.

A: $\answer{1/6}$ \quad B: $\answer{1/6}$ \quad C: $\answer{2/3}$

\begin{feedback}
C is pivotal $4$ times, and A and B once each, out of $6$: $\frac16$, $\frac16$, $\frac46=\frac23$.
\end{feedback}
\end{prompt}

\begin{prompt}
\textbf{(b)} Using the pivots shown, which coalitions are winning?

\begin{selectAll}
\choice{$\set{A,B}$}
\choice[correct]{$\set{A,C}$}
\choice[correct]{$\set{B,C}$}
\choice[correct]{$\set{A,B,C}$}
\choice{$\set{C}$}
\end{selectAll}

\begin{hint}
A pivot at position $k$ means the first $k$ players win together, but the first $k-1$
do not.
\end{hint}

\begin{feedback}
$\seq{A,B,\underline{C}}$ shows $\set{A,B}$ loses and $\set{A,B,C}$ wins.
$\seq{A,\underline{C},B}$ shows $\set{A,C}$ wins, and $\seq{B,\underline{C},A}$ shows
$\set{B,C}$ wins. $\seq{C,\underline{A},B}$ shows $\set{C}$ alone loses. So the winning
coalitions are $\set{A,C}$, $\set{B,C}$, and $\set{A,B,C}$.
\end{feedback}
\end{prompt}

\begin{prompt}
\textbf{(c)} Determine the Banzhaf power distribution, as percentages.

A: $\answer{20}\%$ \quad B: $\answer{20}\%$ \quad C: $\answer{60}\%$

\begin{feedback}
Critical players: both in $\set{A,C}$ and in $\set{B,C}$; only C in $\set{A,B,C}$ (removing
A or B still leaves a winning pair). Counts A $1$, B $1$, C $3$, total $5$: $20\%$,
$20\%$, $60\%$. Note that the two power indices disagree here.
\end{feedback}
\end{prompt}
\end{problem}

\begin{problem}
Consider a city council with a mayor and four council members, $M$, $C_1$, $C_2$, $C_3$,
$C_4$. Everyone has one vote, and a motion passes by simple majority \emph{or} by the mayor
with any one council member's support. Use $*$ for any one council member.

\begin{prompt}
\textbf{(a)} There are five sequential coalition types, depending on where the mayor
falls. In how many of them is the mayor pivotal?

$\answer{2}$

\begin{feedback}
\[
\begin{array}{ll}
\seq{M,\underline{*},*,*,*} & \text{$M$ plus one member wins at step 2}\\
\seq{*,\underline{M},*,*,*} & \text{$M$ joins one member}\\
\seq{*,*,\underline{M},*,*} & \text{two members lose; adding $M$ wins}\\
\seq{*,*,\underline{*},M,*} & \text{three members are a majority}\\
\seq{*,*,\underline{*},*,M} & \text{same}
\end{array}
\]
The mayor is pivotal in $2$ types.
\end{feedback}
\end{prompt}

\begin{prompt}
\textbf{(b)} How many distinct sequential coalitions are there of each type, and in
total?

Each type: $\answer{24}$ \qquad Total: $\answer{120}$

\begin{feedback}
Within a type, the four council members can be arranged in $4!=24$ ways. Five types
give $5\cdot 24=120=5!$. \checkmark
\end{feedback}
\end{prompt}

\begin{prompt}
\textbf{(c)} Determine the Shapley-Shubik power distribution, as percentages.

Mayor: $\answer{40}\%$ \quad Each council member: $\answer{15}\%$

\begin{feedback}
The mayor is pivotal in $2\cdot 24=48$ sequential coalitions: $\frac{48}{120}=40\%$.
Council members are pivotal in the other $72$, split four ways: $18$ each, or $15\%$.
The mayor's one vote carries as much power as almost three council votes.
\end{feedback}
\end{prompt}
\end{problem}

\end{document}
